\section{方法}
\label{sec:method-zh}

DVCL 的核心思想是同时利用两类互补信息。
第一类是经过净化的异构语义拓扑，它保留元路径带来的高阶语义表达能力。
第二类是由节点特征相似性构建的特征诱导视图，它在结构被攻击时提供相对独立的参考。
两个视图分别编码后，通过融合模块得到最终节点表示，并通过跨视图对比学习约束两种表示的一致性。

\subsection{总体框架}

模型输入为受攻击异构图、节点特征、目标类型节点标签以及一组元路径。
DVCL 首先为每条元路径构造语义图，并估计候选边的可靠性；随后进行元路径级过滤和全局语义过滤，得到净化拓扑视图。
与此同时，模型基于目标节点特征相似性构造 KNN 图，形成特征诱导视图。
两个视图经过编码器得到拓扑表示和特征表示，再通过融合模块用于节点分类。

\subsection{净化拓扑视图}

拓扑视图的目标是在保留异构语义结构的同时削弱攻击边影响。
论文中的净化策略包含两个层次：先在每条元路径对应的语义图内估计边置信度并过滤低置信边，再根据语义注意力评估不同元路径的重要性，进行全局语义过滤。
这样可以同时处理局部边噪声和不同语义视图之间的污染程度差异。

\subsection{特征诱导视图}

特征诱导视图从目标节点特征出发，根据节点间相似性构建 KNN 图。
该视图不直接依赖受攻击拓扑，因此能够在结构扰动较强时提供补充信号。
如果两个节点在特征空间中高度相似，即使它们在被攻击图中没有可靠连接，也可以通过该视图建立表示学习联系。

\subsection{双视图编码、融合与分类}

拓扑视图编码器负责从净化后的语义拓扑中学习结构语义表示，特征视图编码器负责从特征诱导图中学习特征相似性表示。
随后，模型将两类表示融合为最终节点表示，并输入分类器完成目标类型节点分类。
融合的目的不是简单替代拓扑信息，而是在拓扑语义和特征稳定性之间取得平衡。

\subsection{跨视图对比学习}

DVCL 使用对称跨视图对比目标对齐两个视图。
对于同一节点，其拓扑表示和特征表示构成正样本对；对于不同节点，表示之间构成负样本对。
这种约束鼓励同一节点在两种视图下保持一致，同时避免所有节点表示塌缩到相同空间位置。

\subsection{训练目标}

最终训练目标由三部分组成：监督节点分类损失、跨视图对比损失和辅助语义拓扑学习损失。
分类损失保证任务性能，对比损失提升双视图一致性，辅助目标帮助拓扑净化过程学习更可靠的语义结构。
